\section*{第 \ref{ch:b3:water-rock-interaction} 章——水–岩相互作用}

\begin{solution}{exo:b3:water-rock-interaction:1}
阳离子：$2\times48/40.08 + 2\times9.7/24.31 + 23/22.99 + 3.9/39.10 =
\qty{4.29}{meq/L}$；阴离子：$183/61.02 + 35.5/35.45 + 2\times48/96.06 =
\qty{5.00}{meq/L}$。误差 $100\times(4.29 - 5.00)/9.29 = \qty{-7.6}{\percent}$：
远超仔细分析此类水样时所允许的 1 或 \qty{2}{\percent}。
可能缺少某种阳离子或其含量被低估，或某种阴离子含量偏高。
\end{solution}

\begin{solution}{exo:b3:water-rock-interaction:2}
$\sqrt I = 0.0707$. \ce{Ca^2+}: $-\log\gamma = 0.5085\times4\times0.0707/(1 +
0.3281\times6\times0.0707) = 0.126$, $\gamma = 0.748$; \ce{HCO3-}: $\gamma =
0.927$; \ce{SO4^2-}: $\gamma = 0.739$.
\end{solution}

\begin{solution}{exo:b3:water-rock-interaction:3}
$\sqrt I = 0.126$; $\gamma_{\ce{Ca}} = 0.622$, $\gamma_{\ce{SO4}} = 0.602$;
$\log\mathrm{IAP} = \log(3.0\times10^{-3}\times0.622\times4.0\times10^{-3}
\times0.602) = -5.35$；$\mathrm{SI} = -5.35 + 4.55 = -0.80$。该水样
未饱和，仍可溶解石膏。
\end{solution}

\begin{solution}{exo:b3:water-rock-interaction:4}
$m = 0.012/60.08 = \qty{2.00e-4}{mol/L}$，$\log m = -3.70$。石英：
$\mathrm{SI} = -3.70 + 3.98 = +0.28$，略过饱和；无定形
二氧化硅：$-3.70 + 2.71 = -0.99$，未饱和。
\end{solution}

\begin{solution}{exo:b3:water-rock-interaction:5}
在 \qty{25}{\celsius} 时，$\log K = -8.45 - 6.35 - 1.47 + 10.33 = -5.94$，$m =
(10^{-5.94}\times10^{-2}/4)^{1/3} = \qty{1.42}{mmol/L}$，\qty{57}{mg/L}；在
\qty{10}{\celsius} 时，$\log K = -5.59$，$m = \qty{1.86}{mmol/L}$，
\qty{75}{mg/L}：低温下约多出三分之一。
\end{solution}

\begin{solution}{exo:b3:water-rock-interaction:6}
$[\ce{CO2}] = 10^{-7.60}\times4.0\times10^{-3}/10^{-6.35} =
\qty{2.25e-4}{mol/L}$; $[\ce{CO3^2-}] = 10^{-10.33}\times4.0\times10^{-3}/
10^{-7.60} = \qty{7.4e-6}{mol/L}$; $\log P_{\ce{CO2}} = -3.648 + 1.47 =
-2.18$，为空气的 16 倍。总量 \qty{4.23}{mmol/L}：\ce{CO2}
\qty{5.3}{\percent}, \ce{HCO3-} \qty{94.5}{\percent}, \ce{CO3^2-}
\qty{0.18}{\percent}.
\end{solution}

\begin{solution}{exo:b3:water-rock-interaction:7}
$r = 10^{-12.09} = \qty{8.2e-13}{mol.m^{-2}.s^{-1}}$（酸性项、中性项和
碱性项分别贡献 $10^{-12.45}$、$10^{-12.56}$ 和 $10^{-12.74}$）。
半径缩小的速率为 $rV_m = 8.2\times10^{-13}\times1.0131\times10^{-4} =
\qty{8.3e-17}{m/s}$，即每年 \qty{2.6}{nm}；半径为 \qty{0.5}{mm} 的颗粒
可维持 $6.0\times10^{12}$~s，约 \num{190000} 年。
\end{solution}

\begin{solution}{exo:b3:water-rock-interaction:8}
$\exp[(69.8/0.008314)(1/278.15 - 1/298.15)] = 7.6$；
若取 \qty{74}{kJ/mol} 则为 8.6。
\end{solution}

\begin{solution}{exo:b3:water-rock-interaction:9}
令 $(12.703 - X - 3Y)/3$ 与 $(7.435 - 2Y)/2$ 相等，消去 $Y$：$X =
12.703 - 11.153 = 1.55$，为一条水平线。该水样的 $X = \log(2.0\times
10^{-3}/39.10) + 7.0 = 2.71$，$Y = \log(0.015/60.08) = -3.60$，该点位于 1.55 之上：
且远低于白云母–钾长石线（$X = -3.161 - 3Y = 7.65$），
故形成白云母。
\end{solution}

\begin{solution}{exo:b3:water-rock-interaction:10}
$\delta_w^f = -2 - 2 = -4$。封闭系统：$W/R = 0.5\times(-2 - 6)/(-13 + 4) =
0.44$。开口系统：$0.5\ln[(-8)/(-9) + 1] = 0.5\ln1.889 = 0.32$。在开口系统中，每个
水团离开时与其最后接触的岩石处于平衡，尽可能多地
带走氧-18；在封闭系统中，整批水最终为 $-4$，
且每千克带走的量更少。
\end{solution}

\begin{solution}{exo:b3:water-rock-interaction:11}
开口系统：$\log K = -5.59$，$m = (10^{-5.59}\times10^{-1.5}/4)^{1/3} =
\qty{2.73}{mmol/L}$，\qty{109}{mg/L}。封闭系统：$C_0 = 10^{-1.27 - 1.5} =
\qty{1.70e-3}{mol/L}$，且 $K_{\mathrm{cal}}K_1/K_2 = 10^{-4.32}$，由三次方程
$4m^3 = 4.79\times10^{-5}(1.70\times10^{-3} - m)$ 可得 $m =
\qty{1.45}{mmol/L}$，\qty{58}{mg/L}，约为原来的一半。剩余的二氧化碳
为 $0.25\times10^{-3}$，占 $C_0$ 的 \qty{15}{\percent}，故 $\log P_{\ce{CO2}}$
从 $-1.5$ 降至 $-2.33$。
\end{solution}

\begin{solution}{exo:b3:water-rock-interaction:12}
$1 - 10^{\mathrm{SI}}$：0.99、0.90 和 0.21。幂律给出
$(10^4/6)^{-0.61} = 0.011$：速率约低 90 倍。野外矿物
已老化，反应性表面已耗尽，并浸泡在
接近平衡的缓慢孔隙水中；两者相乘，野外速率
比实验室速率低若干个数量级。
\end{solution}

\begin{solution}{pb:b3:water-rock-interaction:1}
\textbf{1.} \ce{Ca} 2.795, \ce{Mg} 0.165, \ce{Na} 0.305, \ce{K} 0.031,
\ce{HCO3} 5.572, \ce{Cl} 0.310, \ce{SO4} 0.146, \ce{NO3} \qty{0.145}{mmol/L}.

\textbf{2.} 阳离子 \qty{6.25}{meq/L}，阴离子 \qty{6.32}{meq/L}：误差
\qty{-0.5}{\percent}，可接受。

\textbf{3.} $I = \tfrac12(4\times2.795 + 4\times0.165 + 0.305 + 0.031 + 5.572
+ 0.310 + 4\times0.146 + 0.145)\times10^{-3} = 0.0094$.

\textbf{4.} $\sqrt I = 0.0969$: $\gamma_{\ce{Ca^2+}} = 0.683$,
$\gamma_{\ce{HCO3-}} = 0.904$, $\gamma_{\ce{CO3^2-}} = 0.676$.

\textbf{5.} 是：$I$ 约为限值（约 0.1）的十分之一。

\textbf{6.} $a_{\ce{HCO3-}} = 5.572\times10^{-3}\times0.904 = 5.04\times
10^{-3}$ ($\log = -2.298$).

\textbf{7.} $a_{\ce{CO2}} = 10^{-7.20}\times5.04\times10^{-3}/10^{-6.46} =
9.17\times10^{-4}$ ($\log = -3.038$).

\textbf{8.} $\log P_{\ce{CO2}} = -3.038 + 1.27 = -1.77$, \qty{0.017}{atm}:
$10^{1.60} = 40$ 倍于空气，属地下水的常见范围。二氧化碳
来自补给区土壤空气，源于根系和有机质的分解。

\textbf{9.} $a_{\ce{CO3^2-}} = 10^{-10.49}\times5.04\times10^{-3}/10^{-7.20}
= 2.58\times10^{-6}$，$[\ce{CO3^2-}] = \qty{3.8e-6}{mol/L}$。总量
\qty{6.49}{mmol/L}：\ce{CO2} \qty{14.1}{\percent}，\ce{HCO3-}
\qty{85.8}{\percent}，\ce{CO3^2-} \qty{0.06}{\percent}。

\textbf{10.} 碱度 $5.572 + 2\times0.0038 = \qty{5.580}{meq/L}$：
\qty{99.9}{\percent} 为碳酸氢根。

\textbf{11.} $a_{\ce{Ca^2+}} = 2.795\times10^{-3}\times0.683 = 1.909\times
10^{-3}$; $\log\mathrm{IAP} = \log(1.909\times10^{-3}\times2.58\times10^{-6})
= -8.307$; $\mathrm{SI} = -8.307 + 8.35 = +0.04$.

\textbf{12.} 在 $\pm0.1$ 之内，属于指数的不确定度范围，水
与方解石处于平衡。水在岩石中停留了很长时间，
是一处缓慢、弥散的泉。

\textbf{13.} $m = [10^{-5.59}\times0.0171/(4\times0.683\times0.904^2)]^{1/3}
= \qty{2.70}{mmol/L}$，\qty{108}{mg/L}，与实测的 \qty{2.79}{mmol/L} 相比：
多出 \qty{3}{\percent}，与指数 $+0.04$ 所表示的相符。

\textbf{14.} $a_{\ce{CO2}} = 10^{-1.27 - 3.37} = 2.29\times10^{-5}$;
$a_{\ce{H+}} = 10^{-6.46}\times2.29\times10^{-5}/5.04\times10^{-3}$, pH 8.80;
$a_{\ce{CO3^2-}} = 1.03\times10^{-4}$; $\mathrm{SI} = \log(1.909\times
10^{-3}\times1.03\times10^{-4}) + 8.35 = +1.65$（采用图中完整的
碱度平衡时为 1.61）。

\textbf{15.} 迭代 $m = [10^{-5.59 - 3.37}/(4\gamma_{\ce{Ca}}
\gamma_{\ce{HCO3}}^2)]^{1/3}$，取 $I = 3m$：$m = \qty{0.72}{mmol/L}$
（\qty{29}{mg/L}），$I = 0.0022$。方解石沉淀：$2.795 - 0.718 =
\qty{2.08}{mmol/L}$，即每升约 \qty{210}{mg}。

\textbf{16.} $0.208\times2.0 = \qty{0.42}{g/s}$，\qty{13}{t} 每年，
$13.1/2.71 = \qty{4.8}{m^3}$；在 \qty{400}{m^2} 上，每年 \qty{12}{mm}，
每天 \qty{0.03}{mm}：远低于本章所引冷泉钙华的数值。

\textbf{17.} $X = \log(1.2\times10^{-3}/39.10) + 5.8 = 1.29$; $Y =
\log(9.0\times10^{-3}/60.08) = -3.82$.

\textbf{18.} $X < 1.55$ 且 $Y > -4.39$：高岭石。石英：$\mathrm{SI} =
-3.82 + 3.98 = +0.16$.

\textbf{19.} 低于 $\log a_{\ce{H4SiO4}} = -4.39$ 时，
二氧化硅含量为 \qty{2.4}{mg/L}。

\textbf{20.} 若 $X = 1.55$，则需 pH $= 1.55 + 4.51 = 6.06$。

\textbf{21.} 酸性 $-10.06 - 0.500\times5.8 = -12.96$；中性 $-12.41$；碱性
$-21.20 + 0.823\times5.8 = -16.43$。合计 $\log r = -12.30$，中性项
占主导。

\textbf{22.} \omterm{def:b3:water-rock-interaction:regimes}{供给限制型}：土层厚、坡度缓，意味着侵蚀缓慢；
新鲜长石抵达风化前缘后被充分淋滤，
通量取决于侵蚀速率，而不是气候。

\textbf{23.} $\exp[(74/0.008314)(1/283.15 - 1/298.15)] = 4.9$.

\textbf{24.} 泉水在出口处的方解石\omterm{def:b3:water-rock-interaction:si}{饱和指数}
为 $\mathrm{SI} = \textbf{+0.04}$。
\end{solution}
